\scanpage{25}
% L1–3
l'ordre induit par $(D_0,\prec)$, \deleted{ce qui} correspond à une \deleted{\unsure{application}} \inserted{morphisme} $\sigma:\unsure{\mathbb{\Delta}|\Delta}_1\unsure{\mathrel{-}|\longrightarrow}\mathbb{\Delta}_n$, \unsure{tel|telle} que, posant $z_1=\sigma^*(u)$, on ait $x_0=\delta_0(z_1),y_0=\delta_0(z_1)$, i.e. on
% L4–7
a $x_0\prec_*y_0$. Ainsi l'ordre $\prec$ donné à l'avance sur $D_0$ (satisfaisant \textcircled{b}) se récupère à partir de la donnée quasi-simpliciale $D_*$ qu'il définit, comme étant $\prec_*$.

% L8–11
\uline{Rmq} Supposons donné l'ens ordonné $D_*$, \inserted{satisfaisant \textcircled{a},} et supposons qu'il existe un ordre $\omega$ sur $D_0$ satisfaisant \textcircled{b}. Alors l'ordre opposé $\omega^{\circ}$ satisfait également \textcircled{b}, mais donne lieu bien
% L12–14
sûr à une autre loi quasi-simpliciale sur la famille des $D_n=\{d\in D_*\mid\mathrm{Card}\,I_d=n+1\}$ (sauf si $D_1=0$ i.e. l'ordre de $D_*$ est discret).
% L15–17
Néanmoins, comme les \inserted{réalisations top. des} deux ens. quasi-simpliciaux obtenus sont \unsure{can.} homéom. à la réalisation top \inserted{$|D_*|$} de l'ens ordonné $D_*$,
% L18–20
(qui, elle, ne dépend pas du choix d'un ordre sur $D_0$) on trouve \deleted{que les 2} \inserted{un homéom. can} \deleted{ens quasi-simpliciaux}
\[
|D_*^{\omega}|\simeq|D_*^{\omega'}|
\]
% archive sheet number: 24
